% Options for packages loaded elsewhere
% Options for packages loaded elsewhere
\PassOptionsToPackage{unicode}{hyperref}
\PassOptionsToPackage{hyphens}{url}
\PassOptionsToPackage{dvipsnames,svgnames,x11names}{xcolor}
%
\documentclass[
  letterpaper,
  DIV=11,
  numbers=noendperiod]{scrartcl}
\usepackage{xcolor}
\usepackage{amsmath,amssymb}
\setcounter{secnumdepth}{-\maxdimen} % remove section numbering
\usepackage{iftex}
\ifPDFTeX
  \usepackage[T1]{fontenc}
  \usepackage[utf8]{inputenc}
  \usepackage{textcomp} % provide euro and other symbols
\else % if luatex or xetex
  \usepackage{unicode-math} % this also loads fontspec
  \defaultfontfeatures{Scale=MatchLowercase}
  \defaultfontfeatures[\rmfamily]{Ligatures=TeX,Scale=1}
\fi
\usepackage{lmodern}
\ifPDFTeX\else
  % xetex/luatex font selection
\fi
% Use upquote if available, for straight quotes in verbatim environments
\IfFileExists{upquote.sty}{\usepackage{upquote}}{}
\IfFileExists{microtype.sty}{% use microtype if available
  \usepackage[]{microtype}
  \UseMicrotypeSet[protrusion]{basicmath} % disable protrusion for tt fonts
}{}
\makeatletter
\@ifundefined{KOMAClassName}{% if non-KOMA class
  \IfFileExists{parskip.sty}{%
    \usepackage{parskip}
  }{% else
    \setlength{\parindent}{0pt}
    \setlength{\parskip}{6pt plus 2pt minus 1pt}}
}{% if KOMA class
  \KOMAoptions{parskip=half}}
\makeatother
% Make \paragraph and \subparagraph free-standing
\makeatletter
\ifx\paragraph\undefined\else
  \let\oldparagraph\paragraph
  \renewcommand{\paragraph}{
    \@ifstar
      \xxxParagraphStar
      \xxxParagraphNoStar
  }
  \newcommand{\xxxParagraphStar}[1]{\oldparagraph*{#1}\mbox{}}
  \newcommand{\xxxParagraphNoStar}[1]{\oldparagraph{#1}\mbox{}}
\fi
\ifx\subparagraph\undefined\else
  \let\oldsubparagraph\subparagraph
  \renewcommand{\subparagraph}{
    \@ifstar
      \xxxSubParagraphStar
      \xxxSubParagraphNoStar
  }
  \newcommand{\xxxSubParagraphStar}[1]{\oldsubparagraph*{#1}\mbox{}}
  \newcommand{\xxxSubParagraphNoStar}[1]{\oldsubparagraph{#1}\mbox{}}
\fi
\makeatother


\usepackage{longtable,booktabs,array}
\usepackage{calc} % for calculating minipage widths
% Correct order of tables after \paragraph or \subparagraph
\usepackage{etoolbox}
\makeatletter
\patchcmd\longtable{\par}{\if@noskipsec\mbox{}\fi\par}{}{}
\makeatother
% Allow footnotes in longtable head/foot
\IfFileExists{footnotehyper.sty}{\usepackage{footnotehyper}}{\usepackage{footnote}}
\makesavenoteenv{longtable}
\usepackage{graphicx}
\makeatletter
\newsavebox\pandoc@box
\newcommand*\pandocbounded[1]{% scales image to fit in text height/width
  \sbox\pandoc@box{#1}%
  \Gscale@div\@tempa{\textheight}{\dimexpr\ht\pandoc@box+\dp\pandoc@box\relax}%
  \Gscale@div\@tempb{\linewidth}{\wd\pandoc@box}%
  \ifdim\@tempb\p@<\@tempa\p@\let\@tempa\@tempb\fi% select the smaller of both
  \ifdim\@tempa\p@<\p@\scalebox{\@tempa}{\usebox\pandoc@box}%
  \else\usebox{\pandoc@box}%
  \fi%
}
% Set default figure placement to htbp
\def\fps@figure{htbp}
\makeatother





\setlength{\emergencystretch}{3em} % prevent overfull lines

\providecommand{\tightlist}{%
  \setlength{\itemsep}{0pt}\setlength{\parskip}{0pt}}



 


\KOMAoption{captions}{tableheading}
\makeatletter
\@ifpackageloaded{caption}{}{\usepackage{caption}}
\AtBeginDocument{%
\ifdefined\contentsname
  \renewcommand*\contentsname{Table of contents}
\else
  \newcommand\contentsname{Table of contents}
\fi
\ifdefined\listfigurename
  \renewcommand*\listfigurename{List of Figures}
\else
  \newcommand\listfigurename{List of Figures}
\fi
\ifdefined\listtablename
  \renewcommand*\listtablename{List of Tables}
\else
  \newcommand\listtablename{List of Tables}
\fi
\ifdefined\figurename
  \renewcommand*\figurename{Figure}
\else
  \newcommand\figurename{Figure}
\fi
\ifdefined\tablename
  \renewcommand*\tablename{Table}
\else
  \newcommand\tablename{Table}
\fi
}
\@ifpackageloaded{float}{}{\usepackage{float}}
\floatstyle{ruled}
\@ifundefined{c@chapter}{\newfloat{codelisting}{h}{lop}}{\newfloat{codelisting}{h}{lop}[chapter]}
\floatname{codelisting}{Listing}
\newcommand*\listoflistings{\listof{codelisting}{List of Listings}}
\makeatother
\makeatletter
\makeatother
\makeatletter
\@ifpackageloaded{caption}{}{\usepackage{caption}}
\@ifpackageloaded{subcaption}{}{\usepackage{subcaption}}
\makeatother
\usepackage{bookmark}
\IfFileExists{xurl.sty}{\usepackage{xurl}}{} % add URL line breaks if available
\urlstyle{same}
\hypersetup{
  colorlinks=true,
  linkcolor={blue},
  filecolor={Maroon},
  citecolor={Blue},
  urlcolor={Blue},
  pdfcreator={LaTeX via pandoc}}


\author{}
\date{}
\begin{document}


\section{Proposal for Semester
Project}\label{proposal-for-semester-project}

\textbf{Patterns \& Trends in Environmental Data / Computational
Movement Analysis / Geo 880}

\begin{longtable}[]{@{}
  >{\raggedright\arraybackslash}p{(\linewidth - 2\tabcolsep) * \real{0.2807}}
  >{\raggedright\arraybackslash}p{(\linewidth - 2\tabcolsep) * \real{0.7193}}@{}}
\toprule\noalign{}
\begin{minipage}[b]{\linewidth}\raggedright
Semester:
\end{minipage} & \begin{minipage}[b]{\linewidth}\raggedright
FS26
\end{minipage} \\
\midrule\noalign{}
\endhead
\bottomrule\noalign{}
\endlastfoot
\textbf{Data:} & Individual GPS Trajectories (high frequency) \& POI
(Points of Interest) / Infrastructure Data for Zurich \\
\textbf{Title:} & The 15-Minute City in Practice: Comparing Theoretical
Accessibility and Actual Mobility Behavior in Zurich \\
\textbf{Student 1:} & Lucie Kern \\
\textbf{Student 2:} & Noémie Simon \\
\end{longtable}

\subsection{Abstract}\label{abstract}

This project evaluates the extent to which Zurich aligns with the
``15-minute city'' concept by comparing theoretical accessibility to
actual mobility behavior of an individual in living in Zurich. Using
Google Timeline movement data and accessibility modelling of services
within 15 minutes, the study quantifies mismatches between nearby
available services and daily travel patterns, examining whether mobility
is primarily shaped by infrastructure.

\subsection{Research Questions}\label{research-questions}

\begin{enumerate}
\def\labelenumi{\arabic{enumi}.}
\tightlist
\item
  What proportion of an individual's daily activities in Zurich occurs
  within 15 minutes from home, walking, cycling, or by public transport?
\item
  How does actual mobility behavior differ from the theoretical
  accessibility provided by services and infrastructure?
\item
  Which types of activities (work, entertainment, groceries, sport,
  leisure, social visits) most frequently occur beyond a 15‑minute
  travel time from home?
\end{enumerate}

\subsection{Results / products}\label{results-products}

We will map the 15-minute accessibility areas and the most frequently
visited locations, in addition we will provide quantitative indicators
showing the proportion of activities occurring within or outside
theoretically accessible 15-minute areas, differentiated by transport
mode and activity type. Results will compare theoretical accessibility
with actual destinations, classify local versus non-local activities,
and visualize mobility patterns through heatmaps of visited locations.
We expect to show that, despite Zurich's highly efficient public
transport network, social and professional habits regularly lead
individuals to exceed the 15-minute radius, demonstrating that
accessibility alone does not determine mobility behavior.

\subsection{Data}\label{data}

We will use a dataset of individual movement data collected through
google timeline over 3 months, containing latitude, longitude, timestamp
and transport mode. To ensure privacy, all personally identifiable
information, including the exact home location, will be anonymized and
spatially generalized before analysis. For spatial context, we will
integrate OpenStreetMap (OSM) data via API or geographic extracts,
specifically for Zurich (city). This data will include Points of
Interest (POI): schools, shops, parks, workplace, as well as the public
transport network (VBZ tram stops, ZVV S-Bahn stations). Combining these
sources will allow us to model dynamic isochrones (areas accessible
within 15 min) based on the time of day.

\subsection{Analytical concepts}\label{analytical-concepts}

This project combines concepts from movement analysis and urban
accessibility. Trajectory reconstruction will rebuild the individual's
movements from Google Timeline GPS data, while segmentation methods will
distinguish trips, stops, and activities. Map matching will align GPS
traces with Zurich's street and public transport networks, and travel
mode detection will identify whether movement occurred by walking, tram,
train, or bicycle. Semantic annotation will classify destinations and
activities such as work, shopping, or leisure. The main conceptual
movement spaces are the activity space and the space-time prism. The
activity space represents the areas regularly used by the individual and
will be analysed through trajectories, clustering, and heatmaps to
identify recurring mobility patterns. The space-time prism represents
the theoretical 15-minute accessible area and will be modelled using
accessibility and network analysis based on realistic walking and public
transport travel times. Additional analyses include comparing actual and
theoretical accessibility and calculating a mobility deficit metric
quantifying trips beyond the accessible 15-minute radius.

\subsection{R concepts}\label{r-concepts}

The following R packages will be used: jsonlite, lubridate, tidyverse,
sf, dplyr, ggplot2, and tmap for spatial and mobility analysis. GPS data
from Google Timeline will be converted into sf objects using
st\_as\_sf() and projected with st\_transform() to allow accurate
distance calculations. Functions such as filter(), group\_by(),
mutate(), lead(), lag(), and difftime() will be used to reconstruct
trajectories, calculate travel times, speeds, and identify trips or
stationary activities. st\_distance() will measure distances between
locations and accessibility zones, while segmentation and static point
detection methods will separate stops and movements. Finally,
st\_convex\_hull() will estimate the individual's activity space, and
ggplot2/tmap will visualize trajectories, heatmaps, and accessibility
overlaps.

\subsection{Risk analysis}\label{risk-analysis}

The major challenge lies in the precision and continuity of GPS data:
signal loss or low sampling frequency could skew the calculation of
transport modes and actual travel times. Furthermore, precise modeling
of Zurich's public transport isochrones requires complex schedule data
(GTFS) which can be computationally heavy. Plan B: If real-time public
transport data is too complex to integrate, we will base isochrones on
constant average speeds for trams and buses (e.g., 15 km/h) to
approximate accessibility. If GPS data is too noisy, we will apply
smoothing algorithms or filter our data. Another challenge is to include
all the services that are accessible within the radius. If there is no
satisfactory data on that, a plan B would be to focus on mobility
infrastructure rather than points of interest and replace the 2n
research question by ``How is an individual's observed mobility behavior
in Zurich related to the travel times implied by the local transport and
street network?''

\subsection{Questions?}\label{questions}

\begin{enumerate}
\def\labelenumi{\arabic{enumi}.}
\tightlist
\item
  Our 2nd research question implies we would include all services within
  the 15 minutes radius. Is OpenStreetMap (OSM) appropriate for that and
  for the scope of the project ? Since our transport data includes trips
  outside the Zurich area, would it be necessary to incorporate more
  detailed information for these other visited locations?
\item
  Regarding accessibility modeling: is it valuable to compute it 15 min
  radius per mode of transportation ? or should we simplify ?
\item
  What is the best and most defensible way to classify activities? is it
  by assigning the right category manually ? if yes, how do we document
  it ? or can we apply some sort of logic and algorithm ?
\item
  How can we correct wrong mode of transportation in raw data?
\end{enumerate}




\end{document}
